\section{Introduction}
\label{sec:intro}

Financial asset recommendation (FAR) identifies and ranks financial securities such as stocks, bonds, and mutual funds for investors according to their suitability. It draws on heterogeneous signals: a customer's past transactions and declared risk tolerance on the investor side, and asset prices, returns, and volatility on the market side. FAR therefore sits at the intersection of personalised ranking and regulatory obligation. The Markets in Financial Instruments Directive II (MiFID II) is a foundational European Union regulation that requires firms to assess suitability before recommending a product, so that the product's risk aligns with the customer's declared risk tolerance and loss capacity. A recommender that maximises predicted return or ranking quality without conditioning on this declared risk band can excel on standard benchmarks yet remain a regulatory liability.

FAR-Trans~\cite{sanzcruzado2024fartransinvestmentdatasetfinancial} is the first open FAR benchmark to pair real asset pricing with real retail-investor transactions and a declared MiFID II risk band per customer, collected from a large European financial institution over January 2018 to November 2022. Its accompanying benchmark evaluates FAR algorithms on two axes: ranking accuracy, via normalised Discounted Cumulative Gain (nDCG@10), and realised profitability, via Return on Investment (ROI@10). A subsequent study finds these two axes are themselves negatively correlated on FAR-Trans, so a model that wins on accuracy tends to lose on return~\cite{10.1145/3780097}.

Neither axis measures whether a recommendation is suitable for the customer's risk band. A model can lead the nDCG or ROI leaderboard while routinely proposing assets two or more bands away from the customer's declared tolerance, the precise failure MiFID II exists to prevent. Risk-aware FAR methods address this only after the fact, re-ranking a base model's top list by a risk-tolerance utility~\cite{10.1145/3773966.3779383} rather than conditioning the model on the band. The declared risk band therefore never enters training, and risk suitability is left to a post-hoc filter or ignored outright.

Our objective is to bring risk suitability, the alignment between an asset's risk and the customer's declared band, inside the FAR model and to balance it against return and accuracy rather than trade one for another. We evaluate every model on three axes: return (ROI@10), ranking accuracy (nDCG@10), and risk suitability, which we make measurable by introducing Profile Coherence at $k$ (PC@$k$). PC@$k$ is the share of a customer's top-$k$ list whose asset risk class lies within one band of the customer's declared MiFID II band, on a four-point ordinal scale from Conservative to Aggressive. To capture how well a recommender reconciles the three, we further introduce a single balance score, their harmonic mean, which is high only when return, accuracy, and risk suitability are strong together and falls sharply when any one is sacrificed. PC@$k$ and this balance are themselves new, but our central contribution is a recommender that achieves the balance they measure.

That recommender is risk-based segmentation of LightGCN~\cite{he2020lightgcnsimplifyingpoweringgraph}: we partition customers by declared band and route each into a dedicated sub-model during training, so the band shapes the learned embeddings directly rather than filtering its output. On top of this segmentation baseline we ask where risk suitability is best enforced and study two exclusive interventions: steering the model during training with a coherence margin loss atop Bayesian Personalised Ranking (BPR)~\cite{rendle2012bprbayesianpersonalizedranking} that penalises scoring a discordant asset above a coherent one, or making the data coherent in the first place by regrouping each customer onto the band implied by their revealed purchases.

We make three contributions:
\begin{itemize}
    \item \textbf{An analytical study of risk factors in FAR-Trans.} Of 228{,}241 scoreable Buy transactions, 18.6\% land two or more bands from the customer's declared tolerance, and this discordance is a persistent customer-level trait concentrated at the Conservative and Aggressive extremes rather than transaction-level noise. A transaction-level regression on 208{,}029 priced Buys, with cluster-robust standard errors, shows coherent Buys earn $+2.94$ percentage points higher 6-month realised return than discordant ones, conditional on volatility, segment, and year. This is evidence that risk suitability carries no return penalty.
    \item \textbf{A risk-based segmentation approach for LightGCN.} We route each customer to a per-band sub-model so that the declared risk band conditions the learned embeddings, the baseline methodology on which the experiments build.
    \item \textbf{Experimental studies across the three axes.} Over 69 temporal evaluation splits we audit the strongest FAR-Trans baselines, LightGCN (best accuracy) and Random Forest (best return), and ablate two exclusive interventions on the segmentation baseline: steering the model with a coherence margin loss, and making the data coherent through regrouping. The baselines under-serve specific risk bands despite adequate aggregate scores. The margin loss lifts PC@10 from 0.794 to 0.918 but trades away some accuracy, whereas regrouping lifts PC@10 to 0.949 while also improving nDCG and leaving ROI essentially unchanged, so of the two levers only regrouping raises the harmonic-mean balance of the three axes.
\end{itemize}
